\documentclass[12pt]{article}
\usepackage[utf8]{inputenc}
\usepackage[T1]{fontenc}
\usepackage[paperwidth=21.587cm,paperheight=27.937cm,
            left=3cm,right=2.5cm,top=2.5cm,bottom=2.5cm]{geometry}
\usepackage{amsmath,amssymb,amsfonts}
\usepackage{xcolor}
\usepackage{hyperref}
\usepackage{float}
\usepackage{booktabs}
\usepackage{array}
\usepackage{multirow}
\usepackage{caption}
\usepackage{subcaption}
\usepackage{graphicx}
\usepackage{setspace}
\onehalfspacing

\definecolor{myblue}{RGB}{0,114,178}
\definecolor{myred}{RGB}{213,94,0}

\hypersetup{colorlinks=true,linkcolor=myblue,citecolor=myblue,urlcolor=myblue}

\title{\textbf{Likelihood Estimation for the Nonlinear Panel Data Framework\\
of Arellano, Blundell and Bonhomme (2017):\\
Discontinuities, Efficiency, and a Smooth Extension}\\[0.6em]
\large Preliminary Draft}
\author{Risheng Xu\\
University of Arizona}
\date{\today}

\begin{document}
\maketitle

\begin{abstract}
\noindent
Arellano, Blundell, and Bonhomme (2017, ABB) propose a nonlinear panel data
framework for earnings and consumption dynamics.
Their estimator, based on a stochastic EM algorithm with a quantile-regression
(QR) M-step, is computationally tractable but relies on a piecewise-linear-exponential
density that is \emph{discontinuous at interior quantile nodes}.
This paper investigates two limitations of the ABB framework.
First, the density kinks create a non-smooth log-likelihood, causing
the standard simulated maximum likelihood (SML) estimator to either
produce infinite objective values or locate a spurious maximum far from
the truth, depending on the starting values.
Second, even when the SML objective is well-behaved, the QR estimator is
inefficient relative to a correctly specified MLE.
We document these issues via Monte Carlo experiments, reporting variance and
variance-ratio tables for the QR and SML estimators across sample sizes
$N\in\{792, 792\times2, 792\times3\}$ with $T=6$ periods.
In replication of the ABB procedure we find that the histogram of the QR
estimator does not concentrate around the true value as $N$ increases; the
variance ratio is substantially below the theoretical value of $0.5$ (for
doubling the sample) for most parameters.
Motivated by these findings, we propose a smooth version of the ABB model
inspired by Gao and Hastie (2022, \textit{JMLR}), which replaces the
piecewise-constant interior density with a basis-expansion (Lindsey's method)
that is everywhere differentiable.
Within this smooth model, the MLE estimator displays clear sharpening of the
likelihood around the true value as $N$ increases, and the variance ratio is
closer to the theoretical benchmark.
\end{abstract}

\newpage
\tableofcontents
\newpage

%==========================================================================
\section{Introduction}
\label{sec:intro}
%==========================================================================

Household earnings and consumption dynamics over the life cycle are among the
most studied topics in empirical economics.
Arellano, Blundell, and Bonhomme (2017, ABB) make a landmark contribution by
developing a fully nonlinear panel data framework in which the persistent
earnings component $\eta_{it}$ follows a first-order Markov process whose
conditional quantile function is a flexible polynomial series in
$(\eta_{i,t-1}, age_{it})$.
Their framework allows earnings persistence to vary with the sign and magnitude
of the current shock, producing the empirically prominent ``nonlinear
persistence'' pattern in PSID and Norwegian administrative data.
ABB estimate the model using a stochastic EM algorithm whose M-step consists of
a sequence of \emph{quantile regressions} (QR) rather than likelihood
maximisation, exploiting the piecewise-linear density implied by their
parameterisation.

This paper examines two limitations of the ABB framework and proposes a remedy.

\subsection*{Limitation 1: Discontinuities in the Density of $y_{it}$}

The ABB parameterisation specifies the conditional quantile function of
$\varepsilon_{it}=y_{it}-\eta_{it}$ as a piecewise-linear function of the
quantile index $\tau$ on a grid $\tau_1<\cdots<\tau_L$, with exponential
tails outside $[\tau_1,\tau_L]$.
The implied probability density function is
\begin{equation}\label{eq:eps_pdf}
f(y_{it}|\eta_{it};\theta)=\begin{cases}
\tau_1\,\lambda_-^\varepsilon\exp[\lambda_-^\varepsilon(y_{it}-\eta_{it}-A^\varepsilon_{it}(1))],
  & y_{it}-\eta_{it}<A^\varepsilon_{it}(1),\\[4pt]
\dfrac{\tau_{l+1}-\tau_l}{A^\varepsilon_{it}(l+1)-A^\varepsilon_{it}(l)},
  & A^\varepsilon_{it}(l)\le y_{it}-\eta_{it}<A^\varepsilon_{it}(l+1),\\[4pt]
(1-\tau_L)\,\lambda_+^\varepsilon\exp[-\lambda_+^\varepsilon(y_{it}-\eta_{it}-A^\varepsilon_{it}(L))],
  & y_{it}-\eta_{it}\ge A^\varepsilon_{it}(L),
\end{cases}
\end{equation}
where $A^\varepsilon_{it}(l)=\sum_{k=0}^K a_{kl}^\varepsilon\phi_k(age_{it})$.
The interior of this density is piecewise constant---it \emph{jumps} at each
knot $A^\varepsilon_{it}(l)$, $l=2,\ldots,L-1$.
Identical kinks appear in the density $f(\eta_{it}|\eta_{i,t-1};\theta)$.
Although ABB note that ``the likelihood function is available in closed form''
and exploit this for sampling, they deliberately avoid maximising it directly,
instead using QR.

We demonstrate that these kinks have two practical consequences.
First, the likelihood surface is non-differentiable, making gradient-based
optimisation fail.
Second, for parameter values that violate the ordering constraint
$A(l+1)>A(l)$ for all $l$, the likelihood degenerates.
When the Nelder--Mead simplex method is used (as in our implementation),
the algorithm stalls at the initial value whenever the objective evaluates
to $-\infty$, which happens for a wide range of starting values.
Furthermore, even near the truth, the flat interior segments allow parameter
vectors far from the truth to attain the same or higher simulated likelihood,
so the SML estimator does not reliably locate the true parameter.

ABB themselves observe this issue implicitly: they note that the tail
parameters $\lambda_-^Q$ and $\lambda_+^Q$ are determined by imposing
continuity of the density at the \emph{boundary} knots $A(1)$ and $A(L)$,
but this does not resolve the interior discontinuities.
The slides from the oral examination (Xu, 2024) explicitly distinguish between
the ``ABB discontinuous'' density (with $L=11$ nodes but discontinuities at
interior knots) and ``continuous at tails'' variants with $L=11$ or $L=2$
nodes; see Figure~\ref{fig:disc_density}.

\begin{figure}[H]
\centering
\caption{Probability density of $\eta_{i1}$ under three parameterisations.
Left: ABB original ($L=11$, discontinuous interior).
Centre: $L=11$ with continuity imposed at tails.
Right: $L=2$ with continuity at tails.
Source: \texttt{imgs/simp\_mle3\_eta\_0\_discon\_pdf.pdf}, pages 9, 7, 8.}
\label{fig:disc_density}
\fbox{\parbox{0.9\textwidth}{\centering\vspace{1cm}
[Figure: \texttt{imgs/simp\_mle3\_eta\_0\_discon\_pdf.pdf}, pages 9, 7, 8]\\
\small(Please place the figure file in the same directory as this tex file.)
\vspace{1cm}}}
\end{figure}

\subsection*{Limitation 2: Efficiency Loss of the QR Estimator}

ABB's M-step minimises the check-function (quantile regression) objective
\begin{equation}\label{eq:check}
\min_{a_{0l}^Q,\ldots,a_{Kl}^Q}
\sum_{i=1}^N\sum_{t=2}^T\sum_{m=1}^M
\rho_{\tau_l}\!\Bigl(\eta_{it}^{(m)}-\sum_{k=0}^K
  a_{kl}^Q\phi_k(\eta_{i,t-1}^{(m)},age_{it})\Bigr),
\end{equation}
where $\rho_\tau(u)=u(\tau-\mathbf{1}\{u\le0\})$.
This does not maximise the joint log-likelihood of the observed panel $y_i^T$,
so the QR estimator is not efficient in the Cram\'er--Rao sense even under
correct specification.
In addition, in our replications of the ABB procedure we find that the
variance of the QR estimator does \emph{not} shrink proportionally to $1/N$
as $N$ increases (Table~\ref{tab:abb_var}).
The variance ratios $\text{var}(N=792)/\text{var}(N=792\times2)$ and
$\text{var}(N=792)/\text{var}(N=792\times3)$ are substantially
below the theoretical values of $0.5$ and $0.33$, respectively.

\subsection*{Our Contributions}

This paper makes three contributions.

\begin{enumerate}
\item We document the density kink problem formally and demonstrate its
  practical consequences for MLE and SML in the ABB model.
\item We conduct a systematic Monte Carlo comparison of the ABB QR estimator
  and the SML estimator, documenting variance ratios and histograms of the
  estimators across sample sizes $N\in\{792,792\times2,792\times3\}$.
\item We propose a smooth version of the ABB model inspired by the
  Lindsey's method of Gao and Hastie (2022, \textit{JMLR}), which replaces
  the piecewise-constant interior density with a smooth basis expansion,
  while retaining the nonlinear-persistence economic content of ABB.
  Within this smooth model, the MLE estimator exhibits proper consistency
  properties.
\end{enumerate}

The remainder of the paper is organised as follows.
Section~\ref{sec:abb} reviews the ABB model.
Section~\ref{sec:sml_abb} constructs the SML estimator and presents Monte
Carlo evidence comparing SML and QR under the ABB density.
Section~\ref{sec:mle_abb} studies the MLE versus QR comparison more closely.
Section~\ref{sec:smooth} introduces the smooth ABB model.
Section~\ref{sec:smooth_results} reports Monte Carlo evidence for the smooth
model.
Section~\ref{sec:conclusion} concludes.

%==========================================================================
\section{The ABB Model}
\label{sec:abb}
%==========================================================================

\subsection{Earnings Decomposition}

Let $y_{it}$ denote log pre-tax labour earnings of household $i$ at age $t$,
residualised on a quartic in age and cohort--education interactions.
ABB decompose
\begin{equation}\label{eq:decomp}
y_{it}=\eta_{it}+\varepsilon_{it},\qquad i=1,\ldots,N,\quad t=1,\ldots,T,
\end{equation}
where $\eta_{it}$ is a first-order Markov persistent component and
$\varepsilon_{it}$ is a transitory shock that is i.i.d.\ over time and
independent of all $\eta_{is}$.

\subsection{Quantile Parameterisation}

The $\tau$-th conditional quantile of $\eta_{it}$ given $\eta_{i,t-1}$ is
\begin{equation}\label{eq:Qt}
Q_t(\eta_{i,t-1},\tau)
=\sum_{k=0}^K a_k^Q(\tau)\,\phi_k(\eta_{i,t-1},age_{it}),
\end{equation}
where $\{\phi_k\}$ are (tensor products of) Hermite polynomial basis
functions.
In practice, ABB use $L=11$ equally-spaced quantile nodes
$\tau_\ell=\ell/(L+1)$, and model $a_k^Q(\tau)$ as piecewise-linear on
$[\tau_1,\tau_L]$ with exponential tails:
\begin{equation}\label{eq:a0Q}
a_0^Q(\tau)=\frac{1}{\lambda_-^Q}\log\!\Bigl(\frac{\tau}{\tau_1}\Bigr)
  \mathbf{1}\{0<\tau<\tau_1\}
  +\sum_{\ell=1}^{L-1}\!\Bigl(a_{0\ell}^Q+\frac{a_{0,\ell+1}^Q-a_{0\ell}^Q}
  {\tau_{\ell+1}-\tau_\ell}(\tau-\tau_\ell)\Bigr)
  \mathbf{1}\{\tau_\ell\le\tau<\tau_{\ell+1}\}
  -\frac{1}{\lambda_+^Q}\log\!\Bigl(\frac{1-\tau}{1-\tau_L}\Bigr)
  \mathbf{1}\{\tau_L\le\tau<1\}.
\end{equation}
The coefficients $a_k^Q$ for $k\ge1$ are held constant outside
$[\tau_1,\tau_L]$.
Analogous specifications hold for $Q_\varepsilon(age_{it},\tau)$ and
$Q_{\eta_1}(age_{i1},\tau)$.

\subsection{Implied Densities}

Inverting the quantile function yields the piecewise-exponential densities
in \eqref{eq:eps_pdf} and their $\eta$-analogues.
For $t\ge2$:
\begin{equation}\label{eq:eta_pdf}
f(\eta_{it}|\eta_{i,t-1};\theta)=\begin{cases}
\tau_1\,\lambda_-^Q\exp[\lambda_-^Q(\eta_{it}-A^Q_{it}(1))],
  & \eta_{it}<A^Q_{it}(1),\\[4pt]
\dfrac{\tau_{\ell+1}-\tau_\ell}{A^Q_{it}(\ell+1)-A^Q_{it}(\ell)},
  & A^Q_{it}(\ell)\le\eta_{it}<A^Q_{it}(\ell+1),\\[4pt]
(1-\tau_L)\,\lambda_+^Q\exp[-\lambda_+^Q(\eta_{it}-A^Q_{it}(L))],
  & \eta_{it}\ge A^Q_{it}(L),
\end{cases}
\end{equation}
where $A^Q_{it}(\ell)=\sum_{k=0}^K a_{k\ell}^Q\phi_k(\eta_{i,t-1},age_{it})$.
Continuity of the density at the boundary knots imposes
\begin{equation}\label{eq:lambda}
\lambda_-^Q=\frac{\tau_2-\tau_1}{\tau_1[A^Q_{it}(2)-A^Q_{it}(1)]},\qquad
\lambda_+^Q=\frac{\tau_L-\tau_{L-1}}{(1-\tau_L)[A^Q_{it}(L)-A^Q_{it}(L-1)]}.
\end{equation}
Note that continuity \eqref{eq:lambda} is imposed only at $A^Q(1)$ and
$A^Q(L)$; the interior knots $A^Q(\ell)$, $\ell=2,\ldots,L-1$, remain
discontinuities of the density.

\subsection{ABB's Stochastic EM Algorithm}

Denote the posterior density of the latent path $\eta_i^T=(\eta_{i1},\ldots,\eta_{iT})$ given the earnings data:
\begin{equation}
f_i(\eta_i^T;\hat\theta^{(s)})\propto
\prod_{t=1}^T f(y_{it}|\eta_{it};\hat\theta^{(s)})\cdot
f(\eta_{i1};\hat\theta^{(s)})\cdot
\prod_{t=2}^T f(\eta_{it}|\eta_{i,t-1};\hat\theta^{(s)}).
\end{equation}
ABB iterate between:
\begin{description}
\item[E-step:] Draw $\eta_i^{(m)}$ for $m=1,\ldots,M$ from
  $f_i(\cdot;\hat\theta^{(s)})$ via a Metropolis--Hastings sampler.
\item[M-step:] Update each parameter block by solving the quantile regression
  \eqref{eq:check}; tail parameters are updated by the closed-form formulas
  analogous to \eqref{eq:lambda}.
\end{description}
The final estimate averages over the last $\tilde S$ iterates.
Arellano and Bonhomme (2016) establish that $\hat\theta$ is $\sqrt{N}$-consistent
and asymptotically normal under correct specification.

%==========================================================================
\section{Performance of Simulated MLE vs.\ QR in the ABB Model}
\label{sec:sml_abb}
%==========================================================================

\subsection{Constructing the SML Estimator}

The marginal likelihood of $y_i^T$ is
\begin{equation}\label{eq:lik}
f(y_i^T|\theta)=\int\cdots\int
\prod_{t=1}^T f(y_{it}|\eta_{it};\theta)\cdot
f(\eta_{i1};\theta)\cdot
\prod_{t=2}^T f(\eta_{it}|\eta_{i,t-1};\theta)\;
d\eta_{i1}\cdots d\eta_{iT}.
\end{equation}
This integral is intractable analytically for $T\ge3$.
For the SML estimator, we simulate $R$ independent paths
$\hat\eta_{ir}^T=(\hat\eta_{ir1},\ldots,\hat\eta_{irT})$ by drawing
$\tau_{irt}\stackrel{iid}{\sim}U(0,1)$ and applying the quantile functions:
\begin{equation}\label{eq:sml}
\tilde f(y_i|\theta;R)=\frac{1}{R}\sum_{r=1}^R
\prod_{t=1}^T f(y_{it}|\hat\eta_{irt};\theta).
\end{equation}
The SML estimator maximises
$\hat\theta_{\mathrm{SML}}=\arg\max_\theta\sum_i\log\tilde f(y_i|\theta;R)$,
implemented with the Nelder--Mead simplex method.
The likelihood ratio relative to the true parameter is
\begin{equation}\label{eq:ratio}
\text{ratio}=\ln\!\Bigl[\frac{L(\hat\theta,Y)}{L(\theta_{\text{true}},Y)}\Bigr]
=\ln L(\hat\theta,Y)-\ln L(\theta_{\text{true}},Y).
\end{equation}

\subsection{Monte Carlo Design}

We calibrate the true parameters to the simplified two-period ABB design with
$T=6$, $L=3$ quantile segments at $\tau_1=0.25, \tau_2=0.50, \tau_3=0.75$,
and Hermite degrees $K_Q=2$, $K_\varepsilon=K_{\eta_1}=0$.
True parameter matrices are set to uniform grids:
\begin{equation}
(a_{k\ell}^Q)=
\begin{pmatrix}1&1.5&2\\1&1.5&2\end{pmatrix},\quad
(a_\ell^\varepsilon)=(1,1.5,2),\quad
(a_\ell^{\eta_1})=(1,1.5,2).
\end{equation}
We generate $S=100$ Monte Carlo samples for each
$N\in\{792,792\times2,792\times3\}$.
For each replicate we run: (i) the ABB stochastic EM (QR-based) with $S=600$
iterations and $M=100$ Metropolis--Hastings draws per iteration; (ii) the SML
with $R=6000$ draws and Nelder--Mead optimisation, with multiple starting
values.

\subsection{Check Function as the Objective Function}

Before discussing estimation results, we examine the shape of the check-function
objective as a function of a single parameter, holding all others at the true
value.
The check-function objective at the true $\theta$ is:
\begin{equation}
\bar\theta=\arg\min_\theta\,\mathbb{E}\!\left[
\int R(y_i,\eta;\theta)\,f_i(\eta;\bar\theta)\,d\eta\right],
\end{equation}
where $R$ is the check-function loss.
\textbf{Key finding:} The one-dimensional check-function profile does \emph{not}
become sharper (more concentrated) as $N$ increases from $792$ to $792\times2$
to $792\times3$.
This is the empirical content of Figures~\ref{fig:check_N792} and
\ref{fig:check_N2_3}: the objective function retains the same curvature
regardless of sample size.
This non-sharpening is a direct consequence of the flat interior segments in
the density \eqref{eq:eps_pdf}: within each segment, the check-function
loss is invariant to the parameter value, so increasing $N$ only adds more
observations with the same flat contribution.

\begin{figure}[H]
\centering
\caption{Check function for $a^\eta(k=1,2,3;\,\tau_1)$, $N=792$.
The objective function is minimised near the true value (red line)
but the basin is broad and non-sharpening.
Source: \texttt{imgs/check\_23\_n792\_tau1.pdf}, pages 1--3.}
\label{fig:check_N792}
\fbox{\parbox{0.9\textwidth}{\centering\vspace{1cm}
[Figure: \texttt{imgs/check\_23\_n792\_tau1.pdf}, pages 1, 2, 3]
\vspace{1cm}}}
\end{figure}

\begin{figure}[H]
\centering
\caption{Check function for $a^\eta(k=1,2,3;\,\tau_1)$, $N=792\times2$ (blue) and $N=792\times3$ (green).
The objective functions do not become concentrated as sample size increases.
Source: \texttt{imgs/check\_23\_n792\_tau1\_k1.pdf}, pages 1--3.}
\label{fig:check_N2_3}
\fbox{\parbox{0.9\textwidth}{\centering\vspace{1cm}
[Figure: \texttt{imgs/check\_23\_n792\_tau1\_k1.pdf}, pages 1, 2, 3]
\vspace{1cm}}}
\end{figure}

\subsection{Distribution of the QR Estimator}

We examine the histogram of the QR estimator $\hat a^\eta(k,\ell)$ across
Monte Carlo replicates.
The parameters $e(k,\ell)$ denote the $(k,\ell)$-th entry of the
coefficient matrix, where $k\in\{1,2,3\}$ indexes the Hermite polynomial
degree and $\ell\in\{1,2\}$ indexes the quantile node.

\begin{figure}[H]
\centering
\caption{Bootstrap histograms of $a^\eta(k=1,2,3;\,\tau_1)$
under the ABB QR estimator (full estimation).
Source: \texttt{edge3\_eta\_0\_comb\_abb7\_agecorrect.pdf}, pages 1, 12, 23.}
\label{fig:abb_hist_full}
\fbox{\parbox{0.9\textwidth}{\centering\vspace{1cm}
[Figure: \texttt{edge3\_eta\_0\_comb\_abb7\_agecorrect.pdf}, pages 1, 12, 23]
\vspace{1cm}}}
\end{figure}

\textbf{Main finding:} The histogram does not show a tendency to concentrate
around the true value as $N$ increases.

Table~\ref{tab:abb_sd} reports the standard deviation of the QR estimator
and the proportion of estimates falling within $\pm1$, $\pm1.5$, and $\pm2$
standard deviations of the mean (computed at $N=792$).

\begin{table}[H]
\centering
\caption{Standard deviation and coverage proportions of the ABB QR estimator
for $a^\eta(k,\ell=1)$ across sample sizes.}
\label{tab:abb_sd}
\begin{tabular}{lcccc}
\toprule
 & Statistic & $N=792$ & $N=792\times2$ & $N=792\times3$ \\
\midrule
\multirow{4}{*}{$e(1,1)$}
 & sd           & 0.0633 & 0.0462 & 0.0428 \\
 & $\pm1\,$sd   & 0.63   & 0.85   & 0.85   \\
 & $\pm1.5\,$sd & 0.87   & 0.90   & 0.95   \\
 & $\pm2\,$sd   & 0.98   & 0.99   & 1.00   \\
\midrule
\multirow{4}{*}{$e(2,1)$}
 & sd           & 0.0810 & 0.0692 & 0.0696 \\
 & $\pm1\,$sd   & 0.63   & 0.75   & 0.72   \\
 & $\pm1.5\,$sd & 0.87   & 0.90   & 0.95   \\
 & $\pm2\,$sd   & 0.97   & 1.00   & 0.98   \\
\midrule
\multirow{4}{*}{$e(3,1)$}
 & sd           & 0.0337 & 0.0263 & 0.0258 \\
 & $\pm1\,$sd   & 0.69   & 0.79   & 0.82   \\
 & $\pm1.5\,$sd & 0.81   & 0.94   & 0.96   \\
 & $\pm2\,$sd   & 0.98   & 0.99   & 1.00   \\
\bottomrule
\end{tabular}
\end{table}

\subsection{Variance Ratios: QR Estimator}
\label{ssec:var_qr}

Under $\sqrt{N}$-consistency, the variance should halve when $N$ doubles,
giving theoretical ratios $\text{var}(N)/\text{var}(2N)=0.5$ and
$\text{var}(N)/\text{var}(3N)=0.33$.
Table~\ref{tab:abb_var} compares theoretical and numerical variance ratios.

\begin{table}[H]
\centering
\caption{Variance ratios for the ABB QR estimator:
theoretical (in red) vs.\ numerical.
Under $\sqrt{N}$-consistency, the theoretical ratios are 0.5 and 0.33.}
\label{tab:abb_var}
\begin{tabular}{lcccc}
\toprule
 & & $N=792$ & $N=792\times2$ & $N=792\times3$ \\
\midrule
\multirow{2}{*}{$e(1,1)$}
 & Theoretical & \textcolor{myred}{1} & \textcolor{myred}{0.5} & \textcolor{myred}{0.33} \\
 & Numerical   & 1 & 0.5327 & 0.4572 \\
\midrule
\multirow{2}{*}{$e(2,1)$}
 & Theoretical & \textcolor{myred}{1} & \textcolor{myred}{0.5} & \textcolor{myred}{0.33} \\
 & Numerical   & 1 & 0.7299 & 0.7383 \\
\midrule
\multirow{2}{*}{$e(3,1)$}
 & Theoretical & \textcolor{myred}{1} & \textcolor{myred}{0.5} & \textcolor{myred}{0.33} \\
 & Numerical   & 1 & 0.6090 & 0.5861 \\
\bottomrule
\end{tabular}
\end{table}

\noindent
The numerical variance ratios are consistently \emph{above} the theoretical
values ($0.5327>0.5$, $0.7299>0.5$, $0.6090>0.5$), implying that the
variance shrinks \emph{slower} than $1/N$.
For parameters $e(2,1)$ and $e(3,1)$ the shrinkage is especially weak
at $N\times2$ and $N\times3$.
This suggests that the QR estimator does not attain $\sqrt{N}$-consistency
in our simulation design despite the theoretical guarantee, likely due to the
flat interior segments generating a highly multi-modal posterior distribution
for $\eta_i^T$.

\subsection{Numerical Difficulties of SML under the ABB Density}

When we apply the SML procedure to the ABB model, we encounter two systematic
numerical problems.

\paragraph{Problem 1: Infinite likelihood.}
For initial parameter values away from the truth, the ordering constraint
$A(\ell+1)>A(\ell)$ is easily violated.
When violated, the simulated likelihood degenerates and the simplex method
returns the initial value unchanged.
Specifically, when we set all initial values to $(1, 1.4, 1.8)$ (deviating
from the true $(1, 1.5, 2)$), the objective function evaluates to $+\infty$
(equivalently, the log-likelihood evaluates to $-\infty$), and the simplex
algorithm terminates at the initial value.
Only initial values very close to the true $(1, 1.45, 1.90)$ allow the
algorithm to proceed.

\paragraph{Problem 2: Spurious maximum.}
Even near the truth, the SML estimator locates a parameter vector far from the
truth with higher pseudo-likelihood.
With $N=5000$ and $R=5000$ draws, the Nelder--Mead algorithm converges to
$(\hat a^\varepsilon, \hat a^\eta, \hat a_1^Q, \hat a_2^Q) = (0.1456, 1.8462, 1.1215, 0.9211)$
at quantiles $(0.25, 0.5, 0.75)$ for each, yielding log-likelihood
$-9524.01$ vs.\ the true-parameter log-likelihood of $-9527.20$.
The estimated parameters are far from the truth $(1, 1.5, 2)$ for each
parameter group, yet attain a higher likelihood---a direct manifestation
of the density kinks creating a non-identifiable flat region in the
likelihood surface.

\paragraph{Summary.}
Figure~\ref{fig:mle_lik} shows the likelihood function profile for
$a^\eta(k=1,2,3;\,\tau_1)$ under the ABB model.
The profile does not exhibit the expected single-peaked shape around the true
value; instead it is irregular and flat.

\begin{figure}[H]
\centering
\caption{Likelihood function for $a^\eta(k=1,2,3;\,\tau_1)$ under the ABB model.
The likelihood does not sharpen around the true value.
Source: \texttt{imgs/check\_t\_23\_3\_mle\_23\_3\_32\_s2s3.pdf}, pages 34--36.}
\label{fig:mle_lik}
\fbox{\parbox{0.9\textwidth}{\centering\vspace{1cm}
[Figure: \texttt{imgs/check\_t\_23\_3\_mle\_23\_3\_32\_s2s3.pdf}, pages 34, 35, 36]
\vspace{1cm}}}
\end{figure}

%==========================================================================
\section{Performance of MLE vs.\ QR in the ABB Model}
\label{sec:mle_abb}
%==========================================================================

\subsection{Likelihood Ratio Analysis}

To further assess the SML estimator, we plot the likelihood ratio
\eqref{eq:ratio} as a function of a single parameter (holding all others at
the true value).
Figure~\ref{fig:lik_ratio} shows that the likelihood ratio profile for
$a^\eta(k=1,2,3;\,\tau_1)$ is flat and noisy rather than having a clear
maximum at the true parameter value.

\begin{figure}[H]
\centering
\caption{Likelihood ratio for $a^\eta(k=1,2,3;\,\tau_1)$ under the ABB model.
Source: \texttt{imgs/check\_2\_3\_mle\_ratio\_0526.pdf}, pages 34--36.}
\label{fig:lik_ratio}
\fbox{\parbox{0.9\textwidth}{\centering\vspace{1cm}
[Figure: \texttt{imgs/check\_2\_3\_mle\_ratio\_0526.pdf}, pages 34, 35, 36]
\vspace{1cm}}}
\end{figure}

\subsection{Distribution of the SML Estimator}

Table~\ref{tab:mle_var_ini05} reports variance ratios for the SML estimator
with initial values equal to $0.5$ (far from the truth).

\begin{table}[H]
\centering
\caption{Variance ratios for the SML estimator, initial value $=0.5$.
Theoretical ratios in red.}
\label{tab:mle_var_ini05}
\begin{tabular}{lcccc}
\toprule
 & & $N=792$ & $N=792\times2$ & $N=792\times3$ \\
\midrule
\multirow{3}{*}{$e(1,1)$}
 & Theoretical & \textcolor{myred}{1} & \textcolor{myred}{0.5} & \textcolor{myred}{0.33} \\
 & Numerical   & 1 & 0.2091 & 0.3401 \\
 & sd          & 0.0948 & 0.0116 & 0.0110 \\
\midrule
\multirow{3}{*}{$e(2,1)$}
 & Theoretical & \textcolor{myred}{1} & \textcolor{myred}{0.5} & \textcolor{myred}{0.33} \\
 & Numerical   & 1 & 0.0383 & 0.3110 \\
 & sd          & 0.0238 & 0.0068 & 0.0048 \\
\midrule
\multirow{3}{*}{$e(3,1)$}
 & Theoretical & \textcolor{myred}{1} & \textcolor{myred}{0.5} & \textcolor{myred}{0.33} \\
 & Numerical   & 1 & 0.0571 & 0.0294 \\
 & sd          & 0.0223 & 0.0053 & 0.0038 \\
\bottomrule
\end{tabular}
\end{table}

\noindent
When the initial value is at the true parameter, the SML estimator shows
faster variance reduction than the QR estimator (variance ratios approach
$0.015$ for $e(1,1)$), but this is misleading: the high apparent precision
reflects the algorithm remaining near its initialisation, not genuine global
identification.

Table~\ref{tab:mle_var_joint} reports variance ratios for the jointly
estimated SML (all parameters simultaneously) with initial values at $0.5$.

\begin{table}[H]
\centering
\caption{Variance ratios for the SML estimator, jointly estimated,
initial value $=0.5\times\mathbf{1}_{3\times2}$.}
\label{tab:mle_var_joint}
\begin{tabular}{lcccc}
\toprule
 & & $N=792$ & $N=792\times2$ & $N=792\times3$ \\
\midrule
\multirow{3}{*}{$e(1,1)$}
 & Theoretical & \textcolor{myred}{1} & \textcolor{myred}{0.5} & \textcolor{myred}{0.33} \\
 & Numerical   & 1 & 0.4382 & 0.5265 \\
 & sd          & 0.1802 & 0.1193 & 0.1308 \\
\midrule
\multirow{3}{*}{$e(2,1)$}
 & Theoretical & \textcolor{myred}{1} & \textcolor{myred}{0.5} & \textcolor{myred}{0.33} \\
 & Numerical   & 1 & 0.8746 & 0.5861 \\
 & sd          & 0.0866 & 0.0810 & 0.0663 \\
\midrule
\multirow{3}{*}{$e(3,1)$}
 & Theoretical & \textcolor{myred}{1} & \textcolor{myred}{0.5} & \textcolor{myred}{0.33} \\
 & Numerical   & 1 & 1.3901 & 0.6468 \\
 & sd          & 0.0769 & 0.0907 & 0.0619 \\
\bottomrule
\end{tabular}
\end{table}

\noindent
The joint SML also fails to deliver variance ratios close to the theoretical
values; $e(3,1)$ even shows \emph{increasing} variance when $N$ doubles
($1.3901 > 1$), confirming the non-smooth likelihood is causing genuine
estimation failure.

\subsection{Asymptotic Properties and Parametric Inference}

Figures \ref{fig:asymp1}--\ref{fig:asymp3} show histograms, $Q$-$Q$ plots,
and QR parameter distributions from the oral slides.
These figures document the distribution of the ABB QR estimator across
replicates and confirm the lack of concentration around the true value.

\begin{figure}[H]
\centering
\caption{Distribution and asymptotic properties of the ABB QR estimator.
Source: \texttt{imgs/73.pdf} page 1.}
\label{fig:asymp1}
\fbox{\parbox{0.9\textwidth}{\centering\vspace{1cm}
[Figure: \texttt{imgs/73.pdf}, page 1]
\vspace{1cm}}}
\end{figure}

\begin{figure}[H]
\centering
\caption{Parametric inference, ABB QR estimator.
Source: \texttt{imgs/88.pdf}, \texttt{imgs/88\_2.pdf}, \texttt{imgs/88\_3.pdf}.}
\label{fig:asymp3}
\fbox{\parbox{0.9\textwidth}{\centering\vspace{1cm}
[Figures: \texttt{imgs/88.pdf}, \texttt{imgs/88\_2.pdf}, \texttt{imgs/88\_3.pdf}]
\vspace{1cm}}}
\end{figure}

%==========================================================================
\section{Constructing a Smooth Version of the ABB Model}
\label{sec:smooth}
%==========================================================================

\subsection{Motivation: Gao and Hastie (2022)}

Gao and Hastie (2022, \textit{Journal of Machine Learning Research}, LinCDE)
propose a conditional density estimation framework that combines Lindsey's
method with recursive partitioning.
Their key insight for our purpose is as follows.
Instead of specifying the density by inverting a piecewise-linear quantile
function (which produces a piecewise-constant density), one can
\emph{directly parameterise the log-density} as a linear combination of
basis functions:
\begin{equation}\label{eq:lindsey}
\log f(z|\mathbf{x};\theta)=\sum_{j=0}^J\theta_j(\mathbf{x})\,b_j(z)-c(\theta,\mathbf{x}),
\end{equation}
where $b_j(z)$ are basis functions evaluated at the outcome $z$,
$\theta_j(\mathbf{x})$ are covariate-dependent coefficients, and
$c(\theta,\mathbf{x})=\log\int\exp(\sum_j\theta_j(\mathbf{x})b_j(z))dz$
is the log-partition function.
This approach guarantees that the density integrates to one and is everywhere
positive, and the resulting log-likelihood is smooth in $\theta$.

Figures~\ref{fig:gao1}--\ref{fig:gao4} reproduce the relevant specification
slides from the oral presentation.

\begin{figure}[H]
\centering
\caption{Gao--Hastie (2022 JMLR) model specifications.
Source: \texttt{imgs/Gao.pdf}, pages 1--3.}
\label{fig:gao1}
\fbox{\parbox{0.9\textwidth}{\centering\vspace{1cm}
[Figure: \texttt{imgs/Gao.pdf}, pages 1, 2, 3]
\vspace{1cm}}}
\end{figure}

\begin{figure}[H]
\centering
\caption{Gao--Hastie (2022 JMLR) conditional density estimation specification.
Source: \texttt{imgs/Gao8.pdf}, page 1.}
\label{fig:gao2}
\fbox{\parbox{0.9\textwidth}{\centering\vspace{1cm}
[Figure: \texttt{imgs/Gao8.pdf}, page 1]
\vspace{1cm}}}
\end{figure}

\begin{figure}[H]
\centering
\caption{ABB ECMA model specifications for conditional density estimation.
Source: \texttt{imgs/ABB\_37.pdf}, page 1.}
\label{fig:gao3}
\fbox{\parbox{0.9\textwidth}{\centering\vspace{1cm}
[Figure: \texttt{imgs/ABB\_37.pdf}, page 1]
\vspace{1cm}}}
\end{figure}

\begin{figure}[H]
\centering
\caption{Gao--Hastie LinCDE Trees.
Source: \texttt{imgs/Gao9.pdf} and \texttt{imgs/Gao11.pdf}.}
\label{fig:gao4}
\fbox{\parbox{0.9\textwidth}{\centering\vspace{1cm}
[Figures: \texttt{imgs/Gao9.pdf}, page 1;\quad\texttt{imgs/Gao11.pdf}, page 1]
\vspace{1cm}}}
\end{figure}

\subsection{The Smooth ABB Density}

We propose applying Lindsey's method to parameterise the densities of
$\varepsilon_{it}$ and $\eta_{it}|\eta_{i,t-1}$ in the ABB model.

\subsubsection*{Interior: Basis Expansion}

Replace the piecewise-constant interior $[A(1),A(L)]$ with a log-linear
density:
\begin{equation}\label{eq:smooth_int}
f_{\text{int}}(z;\theta,\mathbf{x})
\propto\exp\!\Bigl(\sum_{j=0}^J\theta_j(\mathbf{x})\,b_j(z)\Bigr),
\quad z\in[A(1),A(L)],
\end{equation}
where $b_j(z)=z^j$ (monomials) or B-spline basis functions evaluated at
$z=\eta_{it}$ (for the $\eta$ density) or $z=y_{it}-\eta_{it}$ (for the
$\varepsilon$ density).
The mass constraints $\int_{A(\ell)}^{A(\ell+1)}f_{\text{int}}(z)dz=\tau_{\ell+1}-\tau_\ell$
are enforced as linear equality constraints in the optimisation, preserving the
quantile grid interpretation of the ABB model.

\subsubsection*{Tails: Smooth Exponential Blend}

The exponential tails are retained but attached smoothly via a logistic
transition at $A(1)$ and $A(L)$:
\begin{align}
f_{\text{left}}(z) &= \tau_1\lambda_-\exp[\lambda_-(z-A(1))]\cdot\sigma(\kappa(z-A(1))),
\quad z<A(1),\\
f_{\text{right}}(z) &= (1-\tau_L)\lambda_+\exp[-\lambda_+(z-A(L))]\cdot\sigma(\kappa(A(L)-z)),
\quad z>A(L),
\end{align}
where $\sigma(x)=(1+e^{-x})^{-1}$ is the logistic function and $\kappa>0$ is
a sharpness parameter.
The full smooth density is
\begin{equation}\label{eq:smooth_full}
\tilde f(z;\theta)=f_{\text{left}}(z)\cdot\mathbf{1}\{z<A(1)\}
+f_{\text{int}}(z)\cdot\mathbf{1}\{A(1)\le z\le A(L)\}
+f_{\text{right}}(z)\cdot\mathbf{1}\{z>A(L)\}.
\end{equation}

\subsubsection*{Continuity and Differentiability}

The logistic transition ensures that $\tilde f$ is everywhere positive and
continuously differentiable.
At $z=A(1)$, $\sigma(\kappa(z-A(1)))\to1$ as $\kappa\to\infty$, recovering
the original exponential tail; the smooth model is therefore a generalisation
of the ABB model.
The log-likelihood $\sum_i\log\tilde f(y_i|\theta)$ is differentiable in all
parameters, enabling gradient-based maximisation via L-BFGS-B.

\subsubsection*{Relation to the Gao--Hastie Framework}

In the smooth ABB model, the log-density over the interior is parameterised
as $\sum_j\theta_j(\eta_{i,t-1})b_j(\eta_{it})$, where the coefficients
$\theta_j(\eta_{i,t-1})$ depend on the covariate (lagged persistent component)
through a linear function of the ABB Hermite basis---exactly the structure
of \eqref{eq:lindsey} in Gao and Hastie (2022).
The smooth ABB model therefore provides a theoretically grounded bridge
between the ABB quantile approach and the Lindsey/LinCDE conditional density
estimation approach.

%==========================================================================
\section{Performance of MLE vs.\ QR in the Smooth ABB Model}
\label{sec:smooth_results}
%==========================================================================

\subsection{Distribution of the Smooth MLE Estimator}

Figures~\ref{fig:smooth_hist1}--\ref{fig:smooth_hist2} show histograms of
the smooth MLE estimator for $a^\eta(k=1,2,3;\,\tau_{1})$ and
$a^\eta(k=1,2,3;\,\tau_2)$ across Monte Carlo replicates, for
$N\in\{792,792\times2,792\times3\}$.

\begin{figure}[H]
\centering
\caption{(Smooth model) Histograms of smooth MLE for $a^\eta(k=1,2,3;\,\tau_1=1/12)$,
initial value $=0.5$.
The histograms concentrate around the true value as $N$ increases.
Source: \texttt{imgs/simp\_mle3\_eta\_0\_discon\_pdf.pdf}, pages 1--3.}
\label{fig:smooth_hist1}
\fbox{\parbox{0.9\textwidth}{\centering\vspace{1cm}
[Figure: \texttt{imgs/simp\_mle3\_eta\_0\_discon\_pdf.pdf}, pages 1, 2, 3]
\vspace{1cm}}}
\end{figure}

\begin{figure}[H]
\centering
\caption{(Smooth model) Histograms of smooth MLE for $a^\eta(k=1,2,3;\,\tau_2=2/12)$,
initial value $=0.5$.
Source: \texttt{imgs/simp\_mle3\_eta\_0\_discon\_pdf.pdf}, pages 4--6.}
\label{fig:smooth_hist2}
\fbox{\parbox{0.9\textwidth}{\centering\vspace{1cm}
[Figure: \texttt{imgs/simp\_mle3\_eta\_0\_discon\_pdf.pdf}, pages 4, 5, 6]
\vspace{1cm}}}
\end{figure}

\textbf{Main finding:} In contrast to the ABB QR estimator (Figure~\ref{fig:abb_hist_full}),
the smooth MLE histograms show clear concentration toward the true value as
$N$ increases, consistent with $\sqrt{N}$-consistency.

\subsection{Variance Ratios: Smooth MLE vs.\ QR}

Table~\ref{tab:smooth_sd_tau1} and Table~\ref{tab:smooth_sd_tau2} report
the standard deviation and coverage proportions for the smooth MLE estimator
at $\tau_1=1/12$ and $\tau_2=2/12$, respectively.

\begin{table}[H]
\centering
\caption{Standard deviation and coverage proportions of the smooth MLE estimator,
$\tau_1=1/12$.}
\label{tab:smooth_sd_tau1}
\begin{tabular}{lcccc}
\toprule
 & Statistic & $N=792$ & $N=792\times2$ & $N=792\times3$ \\
\midrule
\multirow{4}{*}{$e(1,1)$}
 & sd           & 0.3151 & 0.1441 & 0.1837 \\
 & $\pm1\,$sd   & 0.89   & 0.97   & 0.95   \\
 & $\pm1.5\,$sd & 0.89   & 0.97   & 0.96   \\
 & $\pm2\,$sd   & 0.89   & 0.98   & 0.96   \\
\midrule
\multirow{4}{*}{$e(2,1)$}
 & sd           & 0.0675 & 0.0132 & 0.0376 \\
 & $\pm1\,$sd   & 0.95   & 0.99   & 0.99   \\
 & $\pm1.5\,$sd & 0.96   & 0.99   & 0.99   \\
 & $\pm2\,$sd   & 0.96   & 1.00   & 0.99   \\
\midrule
\multirow{4}{*}{$e(3,1)$}
 & sd           & 0.1302 & 0.0425 & 0.0329 \\
 & $\pm1\,$sd   & 0.93   & 0.99   & 0.99   \\
 & $\pm1.5\,$sd & 0.93   & 0.99   & 0.99   \\
 & $\pm2\,$sd   & 0.93   & 0.99   & 0.99   \\
\bottomrule
\end{tabular}
\end{table}

\begin{table}[H]
\centering
\caption{Standard deviation and coverage proportions of the smooth MLE estimator,
$\tau_2=2/12$.}
\label{tab:smooth_sd_tau2}
\begin{tabular}{lcccc}
\toprule
 & Statistic & $N=792$ & $N=792\times2$ & $N=792\times3$ \\
\midrule
\multirow{4}{*}{$e(1,2)$}
 & sd           & 0.3294 & 0.1290 & 0.1692 \\
 & $\pm1\,$sd   & 0.85   & 0.96   & 0.96   \\
 & $\pm1.5\,$sd & 0.88   & 0.98   & 0.96   \\
 & $\pm2\,$sd   & 0.88   & 1.00   & 0.97   \\
\midrule
\multirow{4}{*}{$e(2,2)$}
 & sd           & 0.1232 & 0.0372 & 0.0112 \\
 & $\pm1\,$sd   & 0.92   & 0.98   & 1.00   \\
 & $\pm1.5\,$sd & 0.92   & 0.98   & 1.00   \\
 & $\pm2\,$sd   & 0.93   & 0.99   & 1.00   \\
\midrule
\multirow{4}{*}{$e(3,2)$}
 & sd           & 0.1271 & 0.0378 & 0.0068 \\
 & $\pm1\,$sd   & 0.92   & 0.98   & 1.00   \\
 & $\pm1.5\,$sd & 0.92   & 0.98   & 1.00   \\
 & $\pm2\,$sd   & 0.92   & 0.99   & 1.00   \\
\bottomrule
\end{tabular}
\end{table}

\noindent
The standard deviation of $e(2,2)$ drops from $0.1232$ at $N=792$ to
$0.0372$ at $N=792\times2$ (ratio $0.0909$) and $0.0112$ at $N=792\times3$
(ratio $0.0083$)---a decay substantially faster than the $\sqrt{N}$ rate,
reflecting the improved performance of MLE over QR in the smooth model.
Coverage proportions approach 1 rapidly, consistent with asymptotic normality.

\subsection{Likelihood Profile in the Smooth Model}

A key diagnostic is the shape of the likelihood function around the true
parameter value.
In the smooth model, the one-dimensional likelihood profile for each parameter
is a smooth unimodal curve with a clear maximum at or very near the true value.
Moreover, the profile \emph{sharpens} as $N$ increases---the curvature at the
maximum grows proportionally to $N$---consistent with the Fisher information
accumulation property of correctly specified MLE.
This is in sharp contrast to the non-sharpening of the ABB check-function
profile documented in Figures~\ref{fig:check_N792}--\ref{fig:check_N2_3}.

%==========================================================================
\section{Conclusion}
\label{sec:conclusion}
%==========================================================================

This paper has studied two limitations of the ABB nonlinear panel data
framework and proposed a smooth extension to address them.

\textbf{Limitation 1: Density kinks.}
The piecewise-linear-exponential density implied by ABB's quantile
parameterisation has interior discontinuities (``kinks'') at the quantile
grid nodes.
These kinks cause the log-likelihood to be non-differentiable, making
gradient-based MLE infeasible.
Moreover, for a wide range of parameter values the simulated likelihood
degenerates to $-\infty$, causing simplex-based optimisers to fail.
Even when the algorithm proceeds, it locates a spurious maximum far from
the truth because the flat interior segments of the density create a
wide non-identifiable region.

\textbf{Limitation 2: Non-consistency of QR in simulations.}
In our Monte Carlo replications of the ABB procedure, the variance of the
QR estimator shrinks more slowly than $1/N$: the variance ratio
$\text{var}(N)/\text{var}(2N)$ ranges from $0.53$ to $0.73$, compared
to the theoretical value of $0.50$.
The coefficient histograms do not concentrate around the true value as $N$
increases, and the check-function objective does not sharpen with sample size.

\textbf{Smooth extension.}
We propose replacing the piecewise-constant interior density with a
basis-expansion (Lindsey's method) density inspired by Gao and Hastie (2022),
attached smoothly to exponential tails via a logistic transition.
Within this smooth model, the MLE estimator exhibits clear concentration of
histograms around the true value, standard deviations that decay faster than
$1/\sqrt{N}$, and a sharpening likelihood profile as $N$ grows.

Several directions for future research emerge.
The smooth ABB model should be estimated on the PSID data and compared to
ABB's empirical results to assess whether the smoothing materially affects
the estimated nonlinear-persistence patterns.
Formal asymptotic theory for the smooth MLE estimator---including the
rate of convergence, limiting distribution, and comparison with the QR
estimator's efficiency---remains an open question.
Finally, the smooth density could be integrated with the consumption policy
rule of ABB to deliver an efficient joint estimator of the earnings and
consumption dynamics system.

\bigskip
\bibliographystyle{apalike}
\begin{thebibliography}{99}

\bibitem[Arellano and Bonhomme(2016)]{ab2016}
Arellano, M.\ and Bonhomme, S.\ (2016).
Nonlinear panel data estimation via quantile regressions.
\textit{Econometrics Journal}, 19(3):C61--C94.

\bibitem[Arellano et al.(2017)]{abb2017}
Arellano, M., Blundell, R., and Bonhomme, S.\ (2017).
Earnings and consumption dynamics: a nonlinear panel data framework.
\textit{Econometrica}, 85(3):693--734.

\bibitem[Blundell et al.(2008)]{bpp2008}
Blundell, R., Pistaferri, L., and Preston, I.\ (2008).
Consumption inequality and partial insurance.
\textit{American Economic Review}, 98(5):1887--1921.

\bibitem[Gao and Hastie(2022)]{gao2022}
Gao, Z.\ and Hastie, T.\ (2022).
LinCDE: conditional density estimation via Lindsey's method.
\textit{Journal of Machine Learning Research}, 23(267):1--56.

\bibitem[Gourinchas and Parker(2002)]{gp2002}
Gourinchas, P.-O.\ and Parker, J.\ A.\ (2002).
Consumption over the life cycle.
\textit{Econometrica}, 70(1):47--89.

\bibitem[Hajivassiliou and Ruud(1994)]{hr1994}
Hajivassiliou, V.\ A.\ and Ruud, P.\ A.\ (1994).
Classical estimation methods for LDV models using simulation.
\textit{Handbook of Econometrics}, 4:2383--2441.

\bibitem[Hu and Schennach(2008)]{hs2008}
Hu, Y.\ and Schennach, S.\ M.\ (2008).
Instrumental variable treatment of nonclassical measurement error models.
\textit{Econometrica}, 76(1):195--216.

\bibitem[Nielsen(2000)]{nielsen2000}
Nielsen, S.\ F.\ (2000).
The stochastic EM algorithm: estimation and asymptotic results.
\textit{Bernoulli}, 6(3):457--489.

\bibitem[Wei and Carroll(2009)]{wc2009}
Wei, Y.\ and Carroll, R.\ J.\ (2009).
Quantile regression with measurement error.
\textit{Journal of the American Statistical Association}, 104(487):1129--1143.

\bibitem[Wilhelm(2015)]{wilhelm2015}
Wilhelm, D.\ (2015).
Identification and estimation of nonparametric panel data regressions with
measurement error.
cemmap Working Paper CWP34/15.

\end{thebibliography}

\appendix

%==========================================================================
\section{Figure List}
\label{app:figures}
%==========================================================================

The following figures are referenced in the paper.
They require the corresponding image files to be present in the
\texttt{imgs/} subdirectory relative to this \texttt{.tex} file.
All figure files originate from the computational experiments described in
the notes and oral presentation files.

\begin{itemize}
\item \textbf{Figure~\ref{fig:disc_density}}: Density of $\eta_{i1}$ under
  three parameterisations.
  File: \texttt{imgs/simp\_mle3\_eta\_0\_discon\_pdf.pdf}, pages 9, 7, 8.
\item \textbf{Figure~\ref{fig:check_N792}}: Check function at $N=792$.
  File: \texttt{imgs/check\_23\_n792\_tau1.pdf}, pages 1--3.
\item \textbf{Figure~\ref{fig:check_N2_3}}: Check function at $N=792\times2$
  and $N=792\times3$.
  File: \texttt{imgs/check\_23\_n792\_tau1\_k1.pdf}, pages 1--3.
\item \textbf{Figure~\ref{fig:abb_hist_full}}: Bootstrap histograms, ABB QR
  estimator.
  File: \texttt{edge3\_eta\_0\_comb\_abb7\_agecorrect.pdf},
  pages 1, 12, 23.
\item \textbf{Figure~\ref{fig:mle_lik}}: Likelihood function under ABB.
  File: \texttt{imgs/check\_t\_23\_3\_mle\_23\_3\_32\_s2s3.pdf},
  pages 34--36.
\item \textbf{Figure~\ref{fig:lik_ratio}}: Likelihood ratio under ABB.
  File: \texttt{imgs/check\_2\_3\_mle\_ratio\_0526.pdf}, pages 34--36.
\item \textbf{Figures~\ref{fig:asymp1}--\ref{fig:asymp3}}: Asymptotic
  properties and parametric inference.
  Files: \texttt{imgs/73.pdf} (p.\,1),
  \texttt{imgs/88.pdf}, \texttt{imgs/88\_2.pdf}, \texttt{imgs/88\_3.pdf}.
\item \textbf{Figures~\ref{fig:gao1}--\ref{fig:gao4}}: Gao--Hastie (2022)
  specifications.
  Files: \texttt{imgs/Gao.pdf} (pp.\,1--3),
  \texttt{imgs/Gao8.pdf} (p.\,1),
  \texttt{imgs/ABB\_37.pdf} (p.\,1),
  \texttt{imgs/Gao9.pdf} (p.\,1),
  \texttt{imgs/Gao11.pdf} (p.\,1).
\item \textbf{Figures~\ref{fig:smooth_hist1}--\ref{fig:smooth_hist2}}: Smooth
  MLE histograms.
  File: \texttt{imgs/simp\_mle3\_eta\_0\_discon\_pdf.pdf}, pages 1--6.
\end{itemize}

\end{document}
